\chapter{ThermoDynamics of Big Bang and Big Crunch}\label{ch:bigbang}

\small
\begin{quote}\itshape
The Big Bang is the simplest thermodynamic event in the universe.
A hot, dense, expanding fluid cools. The pressure-work term
$-\,p\,\nabla\!\cdot u$ of \S\ref{sec:expansion_cools} is positive, the
expansion is everywhere, and the cooling is universal.
\end{quote}
\normalsize

\section{Expansion Cools, at Cosmic Scale}

The thermodynamic content of the Big Bang follows directly from the
two primitives of Chapter~\ref{ch:entropy_turbulence}, applied to the
self-gravitating compressible fluid of
Chapter~\ref{ch:cosmology}. The early universe is a hot dense plasma
in which every comoving fluid element has positive divergence:
$\nabla\!\cdot u > 0$ everywhere, at every time. By the
expansion-cools primitive (\S\ref{sec:expansion_cools}), the
internal-energy equation gives
\[
  \dot e \;=\; -\,p\,\nabla\!\cdot u + \mu_{\mathrm{mom}}|\nabla u|^2
  \;<\; 0
\]
in regions where the pressure work dominates the dissipation, i.e.\
where the bulk expansion of the universe is the dominant dynamical
process. Internal energy decreases, temperature drops, and the
universe cools. There is no fine-tuning in this; it is the same
$-p\,\nabla\!\cdot u$ term that cools the gas in the left chamber of
Joule's experiment (Chapter~\ref{ch:joule}), now applied to a gas
whose container is the universe itself.

The cooling timeline is fixed once the equation of state and the
initial conditions are. The radiation-dominated era ($T \gg
10^4\,\mathrm{K}$, $\nabla\!\cdot u \propto t^{-1}$) cools as $T
\propto a^{-1}$ with $a \propto t^{1/2}$; the matter-dominated era
cools as $T \propto a^{-2}$ with $a \propto t^{2/3}$; and the
present-day cosmic microwave background sits at $T \approx 2.7\,\mathrm{K}$,
the remnant of the radiation field after $\sim 13.8$ billion years of
expansion-cooling. No part of this is statistical; all of it is the
expansion-cools primitive applied at cosmic scale and integrated over
cosmic time.

\section{The CMB as Blackbody Cutoff Signature}

The most direct experimental signature of the Big Bang ---  the
cosmic microwave background --- is the radiative half of the story.
At the recombination epoch ($z \approx 1100$, $T \approx 3000\,\mathrm{K}$)
the electrons and protons combined into neutral hydrogen, the
universe became transparent, and the radiation field decoupled from
the matter field. From that moment on, the radiation field expanded
with the universe at fixed photon number, with the spectrum
red-shifted as $\omega \to \omega/a$ and the temperature dropping as
$T \to T/a$. The Planck spectrum maintains its functional form under
this transformation --- a remarkable feature of the blackbody
spectrum, and a non-trivial test of the with-Dynamics reading of
Chapter~\ref{ch:blackbody}: the cutoff at $\omega \sim
k_{\mathrm{B}}T/\hbar$ rescales together with $T$, so a Planck
distribution at $T_{\mathrm{rec}} = 3000\,\mathrm{K}$ at recombination
becomes a Planck distribution at $T_{\mathrm{today}} = 2.7\,\mathrm{K}$
today, observed as the CMB.

The high-frequency tail of the early radiation field --- the part
that, before recombination, was being thermalised by the
free-electron bath at the blackbody cutoff --- has its present-day
counterpart in the small fraction of the CMB energy that even now
continues to be absorbed and re-emitted by intervening matter, with a
cumulative entropy production whose continuum-mechanics realisation
is the cosmological dissipation $D = \int\!\int
\mu_{\mathrm{mom}}|\nabla u|^2\,dx\,dt$ of the gravity-coupled RNS
flow.

\section{The Joule Pattern at Cosmic Scale}

The three-stage Joule pattern that organises this volume applies to
the Big Bang at cosmic scale:

\begin{enumerate}
\item \textbf{Internal energy reservoir.} The initial hot dense plasma
holds an enormous internal energy density. This is the cosmic
analogue of the high-pressure gas in Joule's left chamber.

\item \textbf{Conversion to bulk kinetic energy.} The Hubble expansion
converts internal energy into the bulk kinetic energy of the comoving
fluid through the pressure-work term, exactly as in Joule's free
expansion --- only with the geometry of an isotropic three-dimensional
expansion rather than a channel between two chambers.

\item \textbf{Dissipation into the bath.} The bulk kinetic energy is
in turn dissipated, at every scale at which structure forms, by the
turbulent and shock-front dynamics of the gravity-coupled RNS flow.
This dissipation is what allowed galaxies, stars, and planets to
condense out of an initially smooth distribution --- the
$D$-generating processes of the previous chapter, played out in the
cosmological setting.
\end{enumerate}

The asymmetry between the two halves of the Joule experiment is here
asymmetry between cosmic expansion (bulk, smooth, universal) and
structure formation (local, dissipative, gravitationally-driven). The
universe as a whole is in the expansion-cools half of the cycle. The
local condensations --- galaxy clusters, galaxies, stars, planets ---
are sites where the structure-formation half has run its
$D$-generating dynamics far enough to produce stable, locally bound
configurations of the kind discussed in Chapter~\ref{ch:cosmology}.

\section{Big Crunch: Compression Warms, at Cosmic Scale}

The hypothetical \emph{Big Crunch}, in which a closed universe whose
present expansion is eventually reversed by its own gravity collapses
back to a hot dense state, is the cosmic-scale realisation of the
compression-warms half of the primitive of
\S\ref{sec:expansion_cools}. The pressure-work term that drives the
universal cooling in the expansion-dominated era,
\[
  -\,p\,\nabla\!\cdot u,
\]
reverses sign once $\nabla\!\cdot u$ does. In a contracting universe
the pressure work is positive, internal energy is increasing, and the
universe \emph{warms}. The same mechanism that delivered the
present-day $2.7\,\mathrm{K}$ CMB by 13.8 billion years of
expansion-cooling would, run in reverse over the symmetric
timescale, deliver an arbitrarily hot final state by the same number
of e-foldings of compression-warming.

Big Bang and Big Crunch are therefore not two different physical
mechanisms but the two signs of the same one --- the cosmic-scale
$-p\,\nabla\!\cdot u$ term, with $\nabla\!\cdot u > 0$ in the
expansion era and $\nabla\!\cdot u < 0$ in the contraction era. The
companion \texttt{cosmology\_2d\_cpu.html} simulator can be run in
either mode by choice of initial conditions: with hot dense central
patch and zero initial velocity (Big Bang); or with initially cool
dispersed matter assembled under its own gravity into a single
collapsing mass (Big Crunch). The diagnostics readout shows the
expansion-cooling in the first case and the compression-warming in
the second.

A point worth noting: the irreversible dissipation $D \geq 0$ does
\emph{not} reverse sign with the expansion. Even in a contracting
universe, the local turbulent and shock-front processes that
generate $D$ continue to convert organised kinetic energy into heat.
The Big Crunch, in the with-Dynamics framework, is not a
time-reversal of the Big Bang --- it is the Big Bang's structural
mirror in the bulk pressure-work term, accompanied by the same
unidirectional irreversibility of $D$ in the small-scale dynamics.
What classical state-function thermodynamics would describe by a
``maximum entropy at the moment of maximum expansion'' the
with-Dynamics framework describes by: the integral $\int D\,dt$ over
the whole history continues to grow, even as the bulk pressure work
reverses direction.

\section{Is the Big Bang/Crunch Inherently 3D?}

The expansion-cools / compression-warms mechanism itself is
\emph{dimension-free}. The pressure-work term
$-\,p\,\nabla\!\cdot u$ is a continuum-mechanics object that lives
in any dimension $d$, with the divergence simply summing more terms:
$\nabla\!\cdot u = \sum_{i=1}^d \partial u_i/\partial x_i$. A 1D Big
Bang (a slab expanding along one axis), a 2D Big Bang (a sheet
expanding radially in a plane --- what
\texttt{cosmology\_2d\_cpu.html} simulates), and the 3D Big Bang all
obey the same primitive. The cumulative dissipation
$D = \int\!\int \mu_{\mathrm{mom}}|\nabla u|^2\,dx\,dt$ is also
dimension-free.

But the \emph{specific structure} of our observed universe is
3D-specific:

\begin{itemize}
\item \textbf{Gravity.} Newton's law in $d$ dimensions falls off as
$1/r^{d-2}$. In $d=1$ gravity grows linearly with distance
(collapse is inevitable, no structure formation regime); in $d=2$
it is logarithmic (marginal); only in $d=3$ does it fall off fast
enough to give the Jeans criterion that allows the bulk expansion
to coexist with local condensation into bound structures.

\item \textbf{Stable bound orbits.} Bertrand's theorem singles out
$d=3$ as the dimension in which closed bound orbits exist under an
inverse-square central force. 2D galaxies and stars would form
qualitatively differently; in $d \geq 4$ they would not form stable
bound orbits at all.

\item \textbf{Radiation density of states.} Photon modes per unit
frequency go as $\omega^{d-1}$. The CMB spectrum is a Planck
distribution with $\omega^3$ density of states; in 2D the
corresponding spectrum would have $\omega^1$ scaling, with a
different Stefan--Boltzmann constant and a different
expansion-cooling trajectory for the radiation temperature.

\item \textbf{Jeans criterion.} The condition for gravitational
collapse depends on dimension through both the gravitational
coupling and the pressure response. Only in $d=3$ does the
classical Jeans length come out at the scale that organises
structure formation as we observe it.
\end{itemize}

The 2D simulator \texttt{cosmology\_2d\_cpu.html} is therefore a
qualitative illustration: it captures the bulk pressure-work term
faithfully and demonstrates expansion-cooling and
compression-warming, but does not reproduce the $d=3$-specific
structure-formation regime.

\section{Computational Big Bang and Big Crunch}

For quantitatively faithful Big Bang / Big Crunch simulations the
3D companions are required. The repository ships two
WebGPU-accelerated 3D self-gravitating RNS simulators:

\begin{itemize}
\item \texttt{cosmology\_3d\_gpu\_100.html} --- $100 \times 100 \times
100$ grid ($10^6$ cells), suitable for desktop GPUs.
\item \texttt{cosmology\_3d\_gpu\_200.html} --- $200 \times 200 \times
200$ grid ($8 \times 10^6$ cells), suitable for high-end GPUs.
\end{itemize}

Both implement the same compressible-flow RNS apparatus coupled to a
Poisson-solved gravitational potential, with sliders for the central
patch density and temperature (Big Bang initial condition), for the
gravitational coupling $G$, and for the EOS parameter $\gamma$ that
sets the pressure response. The reader can observe the expansion-cools
phase, the formation of localised gravitational condensations, the
$D$-generating turbulent and shock dynamics that allow those
condensations to virialise into stable bound configurations, and ---
by choice of initial conditions and $G$ --- the compression-warming
of a Big-Crunch-like trajectory.

The Cosmology chapter (Ch.~\ref{ch:cosmology}) develops the
self-gravitating RNS apparatus and uses the simpler
\texttt{cosmology\_2d\_cpu.html} as its working example. The full 3D
companions are referenced there and used here for the explicitly
cosmological Big Bang / Big Crunch initial conditions of the present
chapter.

The cosmological extension of the framework therefore requires no new
physics beyond what Chapters~\ref{ch:entropy_turbulence},
\ref{ch:blackbody}, and~\ref{ch:cosmology} already establish. The
Big Bang, read in the with-Dynamics framework, is the universe as a
single instance of the same compressible-flow apparatus that
organises Joule, refrigerators, heat engines, chemical bonding, and
blackbody radiation. The 2nd Law of the universe is the inevitability
of its dissipation $D$, and the cosmological expansion-cools is the
inevitability of its pressure-work term.
